\chapter{Short-Rate Models}\label{ch:rc:short-rate-models}

Every night a bank values tens of thousands of callable bonds, Bermudan
swaptions and mortgage pools, each an option to exercise on one of many dates
on a whole curve. A model of the swap rate alone cannot do it: exercising at
year three depends on every rate from year three on. The oldest workable answer
models one number, the instantaneous short rate, and derives the whole curve
from it; the version most banks still run in production fits today's curve
exactly, has bond and swaption prices in closed form, and fits on a tree that
values a Bermudan in milliseconds. Its limits are as well known as its uses: one
factor makes every rate move together, and one volatility cannot fit the
swaptions it will be asked to hedge. This chapter builds that model, the
\omterm{def:rc:short-rate-models:hw}{Hull--White model}, fits it to chapter~2's euro curve and to chapter~5's
co-terminal swaptions, checks its tree against its formulas, and measures what
a second factor adds.

\section{The short rate and affine bond prices}

\begin{definition}[Short-rate model]\label{def:rc:short-rate-models:short}
A \emph{short-rate model}\index{short-rate model} specifies the dynamics of the
instantaneous rate $r_t$ under the risk-neutral measure $\mathbb Q$ (One Quant
Book 4, chapter 5), and prices every zero-coupon bond as
$P(t,T) = \E^{\mathbb Q}_t\bigl[\exp(-\int_t^Tr_s\,ds)\bigr]$.
\end{definition}

\begin{definition}[Affine term-structure model]\label{def:rc:short-rate-models:affine}
A \omterm{def:rc:short-rate-models:short}{short-rate model} is an \emph{affine term-structure model}\index{affine
term-structure model} if its bond prices are exponential-affine in the state:
$P(t,T) = \exp\bigl(\mathcal A(t,T)-\mathcal B(t,T)\,x_t\bigr)$ for a state $x_t$ (the
short rate or a vector of factors) and deterministic functions $\mathcal A,
\mathcal B$. By the Feynman--Kac formula (One Quant Book 4, chapter 4) this holds
whenever the drift and the variance of the state are affine in it.
\end{definition}

\begin{definition}[Vasicek model]\label{def:rc:short-rate-models:vasicek}
The \emph{Vasicek model}\index{Vasicek model} (1977) makes the short rate an
Ornstein--Uhlenbeck process, $dr_t = \kappa(\bar x-r_t)\,dt+\sigma\,dW_t$: mean
reversion at speed $\kappa$ to a level $\bar x$, normal increments, negative rates
possible.
\end{definition}

\begin{proposition}[Vasicek bond prices]\label{prop:rc:short-rate-models:vasicek}
With $B(\tau) = (1-e^{-\kappa\tau})/\kappa$,
\[
P(t,T) = \exp\Bigl(\bigl(\bar x-\tfrac{\sigma^2}{2\kappa^2}\bigr)\bigl(B(\tau)-\tau\bigr)
-\tfrac{\sigma^2}{4\kappa}B(\tau)^2 - B(\tau)\,r_t\Bigr),\qquad \tau = T-t,
\]
and the long zero rate tends to $\bar x-\sigma^2/(2\kappa^2)$: convexity pulls long
yields below the mean level.
\end{proposition}
\begin{proof}
$\int_t^Tr_s\,ds$ is normal given $r_t$, with mean $\bar x\tau+(r_t-\bar
x)B(\tau)$ and variance $\frac{\sigma^2}{\kappa^2}\bigl(\tau-B(\tau)-\frac\kappa2
B(\tau)^2\bigr)$; take $\E[e^{-X}] = e^{-\E X+\Var X/2}$.
\end{proof}

A one-factor Gaussian model gives every zero rate $z(t,T) = -\ln P/\tau$ the
normal volatility $\sigma B(\tau)/\tau$. Mean reversion is thus the model's knob
for the \emph{shape} of the volatility curve, not a statement about where rates
are going: with $\kappa=0$ every zero rate has the same volatility, and with
$\kappa>0$ long rates are less volatile than short ones
(\cref{fig:rc:short-rate-models:voltermstructure}).

\begin{omfigure}
\begin{tikzpicture}
\begin{axis}[width=11.5cm,height=5.4cm,
  xlabel={maturity of the zero rate (years)},ylabel={normal volatility (bp)},
  tick label style={font=\small},label style={font=\small},
  xmin=0,xmax=30,ymin=20,ymax=85,grid=major,grid style={gray!25},
  legend columns=3,legend style={font=\footnotesize,at={(0.5,-0.25)},anchor=north,draw=none,fill=none,/tikz/every even column/.append style={column sep=8pt}},
  table/col sep=comma]
\addplot[omDef,thick] table[x=t,y=k0]{figdata/rates-credit-risk/07-short-rate-models/voltermstructure.csv};
\addplot[omThm,very thick] table[x=t,y=k3]{figdata/rates-credit-risk/07-short-rate-models/voltermstructure.csv};
\addplot[omProp,thick,dashed] table[x=t,y=k10]{figdata/rates-credit-risk/07-short-rate-models/voltermstructure.csv};
\legend{$\kappa=0$,$\kappa=3\%$,$\kappa=10\%$}
\end{axis}
\end{tikzpicture}

\omcaption{Normal volatility of zero rates, $\sigma B(\tau)/\tau$, in a one-factor
Gaussian model with $\sigma = 80$ basis points: mean reversion damps the
volatility of long rates (69.1 basis points at ten years with $\kappa=3\%$, 50.6
with $\kappa=10\%$). Data: the chapter's
tutorial.}\label{fig:rc:short-rate-models:voltermstructure}
\end{omfigure}

\section{Vasicek and Hull--White}

Vasicek's model has three parameters and cannot fit today's curve. Hull and
White made the level time-dependent and chose it to fit.

\begin{definition}[Hull--White model]\label{def:rc:short-rate-models:hw}
The \emph{Hull--White model}\index{Hull--White model} (1990) is
$dr_t = (\vartheta(t)-\kappa r_t)\,dt+\sigma(t)\,dW_t$ with $\vartheta(t)$ chosen so
that the model's bond prices equal today's discount factors for every maturity.
Equivalently, $r_t = f(0,t)+x_t$ with
\[
dx_t = \bigl(y(t)-\kappa x_t\bigr)dt+\sigma(t)\,dW_t,\quad x_0=0,\qquad
y(t) = \int_0^te^{-2\kappa(t-u)}\sigma(u)^2\,du .
\]
\end{definition}

\begin{proposition}[Hull--White bonds and bond options]\label{prop:rc:short-rate-models:hwbond}
With $B(t,T) = (1-e^{-\kappa(T-t)})/\kappa$,
\[
P(t,T) = \frac{P(0,T)}{P(0,t)}\exp\Bigl(-B(t,T)\,x_t-\tfrac12B(t,T)^2y(t)\Bigr),
\]
so $\ln P(T,S)$ is normal with standard deviation $\sigma_P = B(T,S)\sqrt{y(T)}$, and a
call at $T$ on the bond maturing at $S$, strike $K$, is worth
$P(0,S)\Phi(h)-KP(0,T)\Phi(h-\sigma_P)$ with $h = \ln\frac{P(0,S)}{KP(0,T)}/\sigma_P+\sigma_P/2$.
\end{proposition}
\admitted\ (a Gaussian computation under the $T$-forward measure; Brigo and
Mercurio, 2006, chapter 3.)

\begin{definition}[Black--Karasinski model]\label{def:rc:short-rate-models:bk}
The \emph{Black--Karasinski model}\index{Black--Karasinski model} (1991) makes the
logarithm of the short rate mean-reverting and normal, $d\ln r_t =
(\vartheta(t)-\kappa\ln r_t)\,dt+\sigma\,dW_t$: rates stay positive, bond prices have no
closed form, and the model is used on trees. It lost ground when rates went
negative.
\end{definition}

\section{Calibration to the curve and to swaptions}

The curve is fitted by construction. The volatility must be fitted to options,
and the options that matter are those the product can be exercised into: for a
ten-year Bermudan callable every year, the \emph{co-terminal} swaptions, one year
into nine, two into eight, up to nine into one.

\begin{definition}[Jamshidian decomposition]\label{def:rc:short-rate-models:jamshidian}
The \emph{Jamshidian decomposition}\index{Jamshidian decomposition} (1989)
writes an option on a coupon bond, in a one-factor model where every bond price
at expiry is a decreasing function of the single state $x_T$, as a portfolio of
options on its zero-coupon bonds: with $x^\star$ the state at which the coupon
bond is worth the strike, each zero-coupon option is struck at that bond's
price at $x^\star$.
\end{definition}

A receiver swaption with strike $K$ is a call at par on a bond paying $K$ each
year and one at maturity, so it is a sum of zero-coupon bond calls priced by
\cref{prop:rc:short-rate-models:hwbond}.

\omcode{code/firm/shortrate/firm_shortrate.py}{73}{89}{European swaptions in Hull--White by Jamshidian's decomposition.}\label{lst:rc:short-rate-models:jamshidian}

\begin{method}[Co-terminal calibration]\label{met:rc:short-rate-models:coterminal}
Fix $\kappa$ (it shapes the volatility term structure; it is often set from
the ratio of volatilities of different tenors, or by desk convention). For the
co-terminal expiries $e_1<\dots<e_m$ into a common final date, take
$\sigma(t)$ piecewise constant on $[e_{i-1},e_i)$ and solve for each level in turn
so that the at-the-money swaption expiring at $e_i$ reprices at its market
volatility, the earlier levels held. Check the resulting $\sigma(t)$ for
smoothness and sign.
\end{method}

\begin{example}[One sigma is not enough]\label{ex:rc:short-rate-models:fit}
On chapter~2's euro OIS curve with $\kappa=3\%$, the co-terminal swaptions into
ten years from chapter~5's cube have at-the-money normal volatilities from
60.8 basis points (one into nine) to 91.0 (nine into one). The best constant
$\sigma$, 86.0 basis points, gives model volatilities between 76.0 and 76.7: it
misses the one-year expiry by 15.3 basis points and the nine-year by 14.2
(\cref{fig:rc:short-rate-models:fit}). A piecewise-constant $\sigma(t)$ fits all
nine exactly, rising from 68.7 to 116.3 basis points before falling to 104.2 in
the last interval (\cref{fig:rc:short-rate-models:sigma}).
\end{example}

\begin{omfigure}
\begin{tikzpicture}
\begin{axis}[width=11.5cm,height=5.4cm,
  xlabel={swaption expiry (years; underlying ends at year ten)},ylabel={ATM normal volatility (bp)},
  tick label style={font=\small},label style={font=\small},
  xmin=0.5,xmax=9.5,ymin=55,ymax=95,grid=major,grid style={gray!25},
  legend columns=3,legend style={font=\footnotesize,at={(0.5,-0.25)},anchor=north,draw=none,fill=none,/tikz/every even column/.append style={column sep=8pt}},
  table/col sep=comma]
\addplot[only marks,mark=*,mark size=2.4pt,black] table[x=e,y=market]{figdata/rates-credit-risk/07-short-rate-models/fit.csv};
\addplot[omDef,thick,mark=square*] table[x=e,y=constant]{figdata/rates-credit-risk/07-short-rate-models/fit.csv};
\addplot[omThm,thick,dashed,mark=triangle*] table[x=e,y=piecewise]{figdata/rates-credit-risk/07-short-rate-models/fit.csv};
\legend{market (cube),Hull--White one $\sigma$,Hull--White $\sigma(t)$}
\end{axis}
\end{tikzpicture}

\omcaption{Co-terminal swaptions into year ten: market volatilities and Hull--White
fits with $\kappa=3\%$. One volatility gives an almost flat line through the
middle; a piecewise-constant volatility reprices every point. Data: chapter~2's
curve and chapter~5's synthetic cube; the chapter's
tutorial.}\label{fig:rc:short-rate-models:fit}
\end{omfigure}

\begin{omfigure}
\begin{tikzpicture}
\begin{axis}[width=11cm,height=5cm,
  xlabel={time (years)},ylabel={$\sigma(t)$ (bp)},
  tick label style={font=\small},label style={font=\small},
  xmin=0,xmax=10,ymin=60,ymax=120,grid=major,grid style={gray!25},
  table/col sep=comma]
\addplot[omThm,very thick,const plot] table[x=t,y=sigma]{figdata/rates-credit-risk/07-short-rate-models/sigma.csv};
\end{axis}
\end{tikzpicture}

\omcaption{The piecewise-constant Hull--White volatility that reprices the nine
co-terminal swaptions. It must rise to make later expiries more volatile than
mean reversion alone would, a sign that $\kappa=3\%$ and the cube's term
structure disagree. Data: the chapter's tutorial.}\label{fig:rc:short-rate-models:sigma}
\end{omfigure}

\section{Trees}

A Bermudan swaption needs the value of continuing at each exercise date, in
every state: backward induction on a lattice. The trinomial tree of One Quant
Book 5, chapter 22, is built here for the Hull--White state.

\begin{method}[The Hull--White tree]\label{met:rc:short-rate-models:tree}
(1) Build a tree for $x$ with $dx = -\kappa x\,dt+\sigma\,dW$: steps $\Delta t$, spacing
$\Delta x = \sigma\sqrt{3\Delta t}$, nodes $j\Delta x$ for $|j|\le j_{\max} =
\lceil0.184/(\kappa\Delta t)\rceil$, three branches from each node with
probabilities that match the mean $-\kappa j\Delta x\Delta t$ and the variance
$\sigma^2\Delta t$, the branching shifted inwards at $\pm j_{\max}$. (2) Shift the tree
by $\alpha_i$ at each step so that it reprices the curve: carry the Arrow--Debreu
prices $Q_{i,j}$ forward and set $\alpha_i = \frac1{\Delta t}\ln\bigl(\sum_j
Q_{i,j}e^{-j\Delta x\Delta t}/P(0,t_{i+1})\bigr)$. The short rate at node $(i,j)$ is
$\alpha_i+j\Delta x$.
\end{method}

\omcode{code/firm/shortrate/firm_shortrate.py}{164}{184}{Fitting the tree to the curve by forward induction of Arrow--Debreu prices.}\label{lst:rc:short-rate-models:tree}

\begin{omfigure}
\begin{tikzpicture}[x=2.3cm,y=0.9cm,font=\small]
\foreach \j in {-1,0,1} {\draw[omDef,thick] (0,0) -- (1,\j);}
\foreach \a in {-1,0,1} {\foreach \j in {-1,0,1} {\draw[omDef,thin] (1,\a) -- (2,\a+\j);}}
\foreach \a in {-2,...,2} {\fill[omThm] (2,\a) circle (2pt);}
\foreach \a in {-1,0,1} {\fill[omThm] (1,\a) circle (2pt);}
\fill[omThm] (0,0) circle (2pt);
\draw[omProp,thick] (2,2) -- (3,2); \draw[omProp,thick] (2,2) -- (3,1); \draw[omProp,thick] (2,2) -- (3,0);
\foreach \a in {0,1,2} {\fill[omProp] (3,\a) circle (2pt);}
\node[anchor=west,font=\footnotesize,align=left] at (3.1,2) {at $j_{\max}$ the branching\\ turns inwards};
\node[below=3pt,font=\footnotesize] at (0,-2.2) {$t_0$};
\node[below=3pt,font=\footnotesize] at (1,-2.2) {$t_1$};
\node[below=3pt,font=\footnotesize] at (2,-2.2) {$t_2$};
\node[below=3pt,font=\footnotesize] at (3,-2.2) {$t_3$};
\node[anchor=east,font=\footnotesize] at (-0.1,0) {$x=0$};
\node[anchor=east,font=\footnotesize] at (1.9,2) {$2\Delta x$};
\end{tikzpicture}

\omcaption{A Hull--White trinomial tree: three branches from each node, spacing
$\Delta x=\sigma\sqrt{3\Delta t}$; once the width reaches $j_{\max}$ the outer nodes
branch inwards, which is how the tree represents mean reversion. The short rate
at each node is its $x$ plus the step's $\alpha_i$.}\label{fig:rc:short-rate-models:tree}
\end{omfigure}

\begin{example}[Tree against formula]\label{ex:rc:short-rate-models:tree}
With $\sigma = 80$ basis points and $\kappa=3\%$, the at-the-money five-into-five payer
(forward 2.933\%, annuity 4.102) is worth 0.026018 per unit notional by
Jamshidian's formula. The tree with monthly steps gives 0.026119, with
fortnightly steps 0.026079, with steps of a week 0.026044 and of half a week
0.026016: the error falls roughly in proportion to the step. The ten-year
discount factor is repriced exactly at every step size, by construction.
\end{example}

\section{Two factors, and what one factor cannot do}

In any one-factor model every rate is a function of the same state, so changes
of all rates are perfectly correlated. Daily changes of the two- and ten-year
Treasury yields had a correlation of 0.769 over 2016--2026 (chapter~6); a
product paying on the difference of two rates (a CMS spread, chapter~6) is
worthless to a one-factor model.

\begin{definition}[Two-factor Gaussian model]\label{def:rc:short-rate-models:g2}
The \emph{two-factor Gaussian model}\index{two-factor Gaussian model} (G2++) sets
$r_t = \varphi(t)+x_t+y_t$ with $dx = -\kappa_1x\,dt+\sigma_1dW^1$, $dy =
-\kappa_2y\,dt+\sigma_2dW^2$, $d\langle W^1,W^2\rangle = \rho\,dt$, and $\varphi$ fitted to
the curve. The correlation of changes of the zero rates of maturities $\tau_1$ and
$\tau_2$ is
\[
\frac{\sum_{a,b}b_a(\tau_1)b_b(\tau_2)c_{ab}}{\sqrt{\sum_{a,b}b_a(\tau_1)b_b(\tau_1)c_{ab}\,\sum_{a,b}b_a(\tau_2)b_b(\tau_2)c_{ab}}},
\quad b_a(\tau)=\frac{B(\kappa_a,\tau)}{\tau},
\]
with $c_{11}=\sigma_1^2$, $c_{22}=\sigma_2^2$, $c_{12}=c_{21}=\rho\sigma_1\sigma_2$.
\end{definition}

\begin{example}[Decorrelating the curve]\label{ex:rc:short-rate-models:g2}
With a slow factor ($\kappa_1=2\%$, $\sigma_1=75$ basis points), a fast one ($\kappa_2=50\%$,
$\sigma_2=90$ basis points) and $\rho=-0.8$, the model's correlation between the two-
and ten-year zero rates is 0.774, close to the Treasury estimate, and between the
two- and thirty-year 0.687 (\cref{fig:rc:short-rate-models:g2}). The strongly
negative $\rho$ is typical: the fast factor moves the front against the slow
factor to produce slope moves.
\end{example}

\begin{omfigure}
\begin{tikzpicture}
\begin{axis}[width=11.5cm,height=5cm,
  xlabel={maturity of the second zero rate (years)},ylabel={correlation with the 2-year},
  tick label style={font=\small},label style={font=\small},
  xmin=0,xmax=30,ymin=0.6,ymax=1.03,grid=major,grid style={gray!25},
  legend columns=2,legend style={font=\footnotesize,at={(0.5,-0.25)},anchor=north,draw=none,fill=none,/tikz/every even column/.append style={column sep=8pt}},
  table/col sep=comma]
\addplot[omThm,very thick] table[x=t,y=g2]{figdata/rates-credit-risk/07-short-rate-models/g2.csv};
\addplot[omDef,thick,dashed] table[x=t,y=one]{figdata/rates-credit-risk/07-short-rate-models/g2.csv};
\legend{two-factor Gaussian model,any one-factor model}
\end{axis}
\end{tikzpicture}

\omcaption{Correlation of instantaneous changes of the two-year zero rate with
zero rates of other maturities: one in any one-factor model, decaying with
maturity in the two-factor model of \cref{ex:rc:short-rate-models:g2}. Data: the
chapter's tutorial.}\label{fig:rc:short-rate-models:g2}
\end{omfigure}

\section{Tutorial: Hull--White from curve to tree}\label{tut:rc:short-rate-models:tutorial}

\begin{tutorial}
\textbf{Goal.} Fit Hull--White to a curve and to co-terminal swaptions, and check
its tree against its closed form. \textbf{End state:}
\cref{fig:rc:short-rate-models:fit,fig:rc:short-rate-models:sigma} and the numbers of
\cref{ex:rc:short-rate-models:tree}.
\begin{tutsteps}
\item \textbf{The model}: \texttt{HullWhite(curve, kappa, sigmas, knots)} reprices
every discount factor of chapter~2's €STR curve.
\item \textbf{Calibrate}: \texttt{calibrations()} fits one $\sigma$ and a
piecewise-constant $\sigma(t)$ to the nine co-terminals;
\texttt{fit\_table()} converts model prices back to normal volatilities.
\item \textbf{Tree}: \texttt{tree\_check(dt=\ldots)} for monthly to half-weekly steps.
\item \textbf{Two factors}: \texttt{g2\_table()}; \texttt{fig\_rc\_shortrate.py}
writes the charts.
\end{tutsteps}
\textbf{What to change next.} Calibrate with $\kappa=0$ and $\kappa=10\%$ and compare
the shapes of $\sigma(t)$; add a time-dependent $\sigma$ to the tree by varying the
step size.
\end{tutorial}

\section{Build: the short-rate engine}\label{bld:rc:short-rate-models:shortrate}

\begin{build}
\bfield{Purpose} The model behind the book's callable products: Bermudans
(chapter~9), mortgages (chapter~12), exposure simulation (chapter~17).
\bfield{Interface} \texttt{HullWhite(curve, kappa, sigmas, knots)} with
\texttt{bond}, \texttt{bond\_vol}, \texttt{zbo}, \texttt{swaption},
\texttt{annuity\_forward}; \texttt{calibrate\_coterminal}; \texttt{HWTree(curve,
kappa, sigma, horizon, dt)} with \texttt{step\_back}, \texttt{zero\_bonds\_at},
\texttt{european\_swaption}; \texttt{vasicek\_bond}, \texttt{zero\_rate\_vol},
\texttt{g2\_correlation}.
\bfield{Rules} The curve is any object with \texttt{df\_t}; swaps with annual unit
accruals in model time; numpy only.
\bfield{Acceptance tests} \texttt{code/firm/shortrate/tests/}: bonds reprice the
curve and bond options satisfy parity; payer minus receiver is the forward
swap; the tree reprices discount factors and converges to Jamshidian's price as
the step shrinks; the piecewise calibration reprices its targets; Vasicek's long
rate is $\bar x-\sigma^2/2\kappa^2$; G2++ with one factor switched off has
correlation one.
\bfield{Stretch} A time-dependent $\sigma$ on the tree; a G2++ swaption formula;
calibration of $\kappa$ to the ratio of co-terminal and column volatilities.
\end{build}

\begin{omsources}
\item O.\ Vasicek, ``An equilibrium characterization of the term structure'',
\emph{Journal of Financial Economics} 5(2), 1977.
\item J.\ Hull and A.\ White, ``Pricing interest-rate-derivative securities'',
\emph{Review of Financial Studies} 3(4), 1990.
\item F.\ Jamshidian, ``An exact bond option formula'', \emph{Journal of Finance}
44, 1989.
\item F.\ Black and P.\ Karasinski, ``Bond and option pricing when short rates are
lognormal'', \emph{Financial Analysts Journal} 47(4), 1991.
\item D.\ Brigo and F.\ Mercurio, \emph{Interest Rate Models: Theory and Practice},
2nd ed., Springer, 2006, chapters 3--4.
\end{omsources}

\section{Exercises}

\begin{exercise}[$\star$]\label{exo:rc:short-rate-models:1}
In a one-factor Gaussian model with $\sigma = 80$ basis points and $\kappa = 3\%$,
what is the normal volatility of the ten-year and of the thirty-year zero rate?
\end{exercise}

\begin{exercise}[$\star$]\label{exo:rc:short-rate-models:2}
Why can Vasicek's model not reprice today's curve, and what does Hull--White
change to fix it?
\end{exercise}

\begin{exercise}[$\star$]\label{exo:rc:short-rate-models:3}
In a \omterm{def:rc:short-rate-models:vasicek}{Vasicek model} with $\bar x = 3\%$, $\sigma = 1\%$ and $\kappa = 10\%$, what is the
long-maturity zero rate?
\end{exercise}

\begin{exercise}[$\star\star$]\label{exo:rc:short-rate-models:4}
For the tree with monthly steps and $\kappa=3\%$, compute $j_{\max}$ and the number of
nodes at the widest step. How does it change with weekly steps?
\end{exercise}

\begin{exercise}[$\star\star$]\label{exo:rc:short-rate-models:5}
Explain why Jamshidian's decomposition fails in a two-factor model.
\end{exercise}

\begin{exercise}[$\star\star$]\label{exo:rc:short-rate-models:6}
Read \cref{fig:rc:short-rate-models:sigma}. What would the calibrated $\sigma(t)$ look
like with $\kappa=10\%$, and why?
\end{exercise}

\begin{exercise}[$\star\star\star$]\label{exo:rc:short-rate-models:7}
\emph{Coding.} With \texttt{tree\_check}, tabulate the tree's error against
Jamshidian's price for monthly, fortnightly, weekly and half-weekly steps, and
estimate the order of convergence.
\end{exercise}

\begin{exercise}[$\star\star\star$]\label{exo:rc:short-rate-models:8}
\emph{Find the flaw.} ``Our one-factor \omterm{def:rc:short-rate-models:hw}{Hull--White model} reprices every co-terminal
swaption, so it prices our \omterm{def:rc:convexity-adjustments-and-constant-maturity-products:spread}{CMS spread options} correctly.'' Correct it.
\end{exercise}

\section{Problem: One Sigma Is Not Enough}

\begin{problem}[{Weekend problem --- calibrating a production Hull--White model}]\label{pb:rc:short-rate-models:1}
A bank's overnight batch prices ten-year euro Bermudans callable every year with
Hull--White, $\kappa=3\%$, on chapter~2's €STR curve. The model validator asks why the
front office uses a piecewise-constant $\sigma(t)$ when the documentation says
``Hull--White with one volatility''.

\medskip\noindent\textbf{Part I --- The market.}
\begin{enumerate}
\item Which swaptions should the model fit, and why those?
\item Give their volatility range and its shape.
\item What does a rising co-terminal volatility say about forward volatility?
\item Why is the curve not a calibration target?
\item Why is $\kappa$ fixed rather than fitted?
\end{enumerate}

\medskip\noindent\textbf{Part II --- One sigma.}
\begin{enumerate}[resume]
\item Give the best constant $\sigma$.
\item Give the model's volatility range.
\item Give the errors at one and nine years.
\item Which exercises of the Bermudan are overpriced, which underpriced?
\item Why can no single $\sigma$ fix it with $\kappa=3\%$?
\end{enumerate}

\medskip\noindent\textbf{Part III --- Piecewise sigma.}
\begin{enumerate}[resume]
\item Give the first and highest levels of $\sigma(t)$.
\item Why does it fall in the last interval?
\item Does the fitted model price the co-terminals exactly? Other swaptions?
\item What does it assume about volatility beyond year nine?
\item What is the risk of fitting nine numbers every day?
\end{enumerate}

\medskip\noindent\textbf{Part IV --- Judgement.}
\begin{enumerate}[resume]
\item What would you write in the model documentation?
\item Which test would convince the validator?
\item What would a two-factor model add for this product, and what would it cost?
\item State the \emph{named result}: the errors with one $\sigma$ and the range of the
calibrated $\sigma(t)$.
\item In one sentence: what does a production \omterm{def:rc:short-rate-models:short}{short-rate model} need to fit?
\end{enumerate}
\end{problem}

\section{Interview questions}

\begin{interviewq}[$\star$ \iqroles{researcher, bank}]\label{iq:rc:short-rate-models:1}
What does mean reversion do in the \omterm{def:rc:short-rate-models:hw}{Hull--White model}, in terms a trader would
use?
\end{interviewq}

\begin{interviewq}[$\star\star$ \iqroles{researcher}]\label{iq:rc:short-rate-models:2}
Derive the zero-coupon bond price in the \omterm{def:rc:short-rate-models:vasicek}{Vasicek model}.
\end{interviewq}

\begin{interviewq}[$\star\star$ \iqroles{researcher, developer}]\label{iq:rc:short-rate-models:3}
How do you fit a trinomial tree to today's \omterm{def:rc:multi-curve-and-collateral-discounting:curves}{discount curve}?
\end{interviewq}

\begin{interviewq}[$\star\star$ \iqroles{trader, researcher}]\label{iq:rc:short-rate-models:4}
Which swaptions do you calibrate a one-factor model to for a ten-year Bermudan,
and why not the whole matrix?
\end{interviewq}

\begin{interviewq}[$\star\star\star$ \iqroles{researcher, risk}]\label{iq:rc:short-rate-models:5}
What can a one-factor model not price, and how do you know when it matters?
\end{interviewq}

\begin{interviewq}[$\star\star\star$ \iqroles{developer}]\label{iq:rc:short-rate-models:6}
Your tree price of a European swaption differs from the closed form by 0.5\%.
How do you find out why?
\end{interviewq}
